% !TEX program = xelatex
\documentclass[final]{beamer}

\usepackage[size=a0,orientation=landscape,scale=1.00]{beamerposter}
\usepackage{fontspec}
\usepackage{amsmath,amssymb,mathtools}
\usepackage{booktabs,array}
\usepackage{ragged2e}
\usepackage{tikz}
\usetikzlibrary{arrows.meta,calc,positioning}
\usepackage[most]{tcolorbox}

\setsansfont{TeX Gyre Heros}
\setmonofont{TeX Gyre Cursor}

\definecolor{Navy}{HTML}{0B1F33}
\definecolor{NavyTwo}{HTML}{163A53}
\definecolor{Teal}{HTML}{0E8F8B}
\definecolor{TealLight}{HTML}{DDF4F1}
\definecolor{Amber}{HTML}{E8A72E}
\definecolor{AmberLight}{HTML}{FFF2D5}
\definecolor{Ink}{HTML}{16212B}
\definecolor{Muted}{HTML}{536573}
\definecolor{Paper}{HTML}{F6F8F9}
\definecolor{White}{HTML}{FFFFFF}
\definecolor{Rule}{HTML}{CBD5DB}
\definecolor{Red}{HTML}{B53A4C}

\setbeamercolor{background canvas}{bg=Paper}
\setbeamercolor{normal text}{fg=Ink,bg=Paper}
\setbeamertemplate{navigation symbols}{}
\setbeamertemplate{headline}{}
\setbeamertemplate{footline}{}
\setbeamersize{text margin left=2.5cm,text margin right=2.5cm}
\hypersetup{
  pdftitle={A precise causal-set question for David A. Meyer},
  pdfauthor={Ignacio Martín Gandul},
  pdfsubject={Expert adjudication of an all-dimensional F1+F2 not-implies-F3 construction},
  pdfkeywords={causal sets, David Meyer, Madsen, Muller, de Sitter, higher dimensions, longest chain, F1, F2, F3}
}

\newtcolorbox{posterbox}[2][]{
  enhanced,
  colback=White,
  colframe=Rule,
  boxrule=0.9pt,
  arc=3mm,
  left=9mm,right=9mm,top=7mm,bottom=8mm,
  before skip=0mm,after skip=10mm,
  title={#2},
  colbacktitle=Navy,
  coltitle=White,
  fonttitle=\bfseries\fontsize{30}{34}\selectfont,
  attach boxed title to top left={xshift=0mm,yshift=-1mm},
  boxed title style={arc=2.5mm,boxrule=0pt,left=8mm,right=8mm,top=3mm,bottom=3mm},
  #1
}

\newtcolorbox{questionbox}[1][]{
  enhanced,
  colback=TealLight,
  colframe=Teal,
  boxrule=2.4pt,
  arc=3mm,
  left=14mm,right=14mm,top=9mm,bottom=10mm,
  before skip=0mm,after skip=10mm,
  #1
}

\newtcolorbox{scopebox}[1][]{
  enhanced,
  colback=AmberLight,
  colframe=Amber,
  boxrule=1.4pt,
  arc=3mm,
  left=9mm,right=9mm,top=7mm,bottom=7mm,
  before skip=0mm,after skip=10mm,
  #1
}

\newcommand{\badge}[2]{%
  \tikz[baseline=(X.base)]\node[
    rounded corners=2.4mm,
    inner xsep=6mm,
    inner ysep=2.8mm,
    fill=#1,
    text=White,
    font=\bfseries\fontsize{16}{19}\selectfont
  ] (X) {#2};%
}

\newcommand{\key}[1]{\textcolor{Teal}{\bfseries #1}}
\newcommand{\warn}[1]{\textcolor{Red}{\bfseries #1}}
\newcommand{\tightdisplay}[1]{\vspace{-1mm}\[#1\]\vspace{-2mm}}

\begin{document}
\begin{frame}[t]

% Header
\begin{tcolorbox}[
  enhanced,colback=Navy,colframe=Navy,boxrule=0pt,arc=0mm,
  left=16mm,right=16mm,top=10mm,bottom=10mm,width=\textwidth
]
\begin{columns}[T,totalwidth=\linewidth]
  \begin{column}{0.82\linewidth}
    {\fontsize{61}{66}\selectfont\bfseries\color{White}
      A precise causal-set question for David A. Meyer\par}
    \vspace{4mm}
    {\fontsize{31}{36}\selectfont\color{TealLight}
      Does a vanishing-mass planted chain give a genuine
      \(\mathrm{F1}+\mathrm{F2}\not\Rightarrow\mathrm{F3}\) counterexample?\par}
    \vspace{5mm}
    {\fontsize{19}{23}\selectfont\color{White}
      Ignacio Martín Gandul \quad\textcolor{Teal}{\textbullet}\quad
      \texttt{nachocausal} research programme \quad\textcolor{Teal}{\textbullet}\quad
      WP7 · 11 August 2026}
  \end{column}
  \begin{column}{0.18\linewidth}
    \raggedleft
    \begin{tikzpicture}
      \fill[Teal] (0,0) circle (2.2cm);
      \draw[White,line width=2.5pt] (0,0) circle (2.2cm);
      \node[White,font=\bfseries\fontsize{36}{39}\selectfont] at (0,0.2) {?};
      \node[White,font=\bfseries\fontsize{15}{18}\selectfont] at (0,-0.75) {WP7};
    \end{tikzpicture}
  \end{column}
\end{columns}
\end{tcolorbox}

\vspace{8mm}
\fontsize{24}{29}\selectfont

% The one question
\begin{questionbox}
  \begin{columns}[T,totalwidth=\linewidth]
    \begin{column}{0.73\linewidth}
      {\fontsize{34}{40}\selectfont\bfseries\color{Navy}
      Under Madsen's regional reading, does the fixed-\(d\) family
      \(P_\rho=\Pi_\rho\cup\Gamma_\rho\), now proved for every \(d\ge2\), count as a
      faithful, count-uniform counterexample
      to \(\mathrm{F1}+\mathrm{F2}\Rightarrow\mathrm{F3}\) --- or is there an intended
      causal-set condition that excludes it?\par}
    \end{column}
    \begin{column}{0.25\linewidth}
      \centering
      \badge{NavyTwo}{MATHEMATICS CLOSED UNDER STATED READING}
      \vspace{5mm}

      \badge{Amber}{EXPERT ADJUDICATION REQUESTED}
      \vspace{5mm}

      \badge{Teal}{PROVED REGIME: EVERY FIXED \(d\ge2\)}
    \end{column}
  \end{columns}
  \vspace{5mm}
  \textbf{We are not asking whether the two estimates can be proved: they can.}
  We are asking whether the construction matches the intended causal-set semantics and whether
  the same mechanism is already present in the literature.
\end{questionbox}

\begin{columns}[T,totalwidth=\textwidth]

% ------------------------------------------------------------------
% COLUMN 1
% ------------------------------------------------------------------
\begin{column}{0.319\textwidth}

\begin{posterbox}{What Madsen separates}
  For a well-conditioned causal-set embedding:
  \vspace{4mm}

  \renewcommand{\arraystretch}{1.30}
  \begin{tabular}{@{}>{\centering\arraybackslash\bfseries\color{Navy}}m{0.10\linewidth}>{\raggedright\arraybackslash}m{0.84\linewidth}@{}}
    F1 & Exact preservation of the induced causal order.\\
    F2 & Uniform number--volume control in every admissible mesoscopic diamond.\\
    F3 & Longest-chain distance approximates Lorentzian proper time at the prescribed rate.\\
  \end{tabular}

  \vspace{5mm}
  With \(\mathcal N_\rho=\rho V_K\), the relevant F2 tolerance is
  \tightdisplay{
    \bigl|\#(P_\rho\cap D)-\rho\operatorname{Vol}_g(D)\bigr|
    \le K_2\sqrt{\rho\operatorname{Vol}_g(D)\log\mathcal N_\rho}.
  }
\end{posterbox}

\begin{posterbox}{The proposed family}
  \begin{columns}[T,totalwidth=\linewidth]
    \begin{column}{0.48\linewidth}
      \begin{center}
      \begin{tikzpicture}[x=0.55cm,y=0.55cm]
        \fill[Paper] (0,0) rectangle (14,14);
        \draw[Navy,line width=1.8pt] (0,0) rectangle (14,14);
        \draw[Rule,line width=0.8pt,step=2] (0,0) grid (14,14);
        \fill[TealLight,opacity=.75] (2.1,1.7) rectangle (11.8,10.7);
        \draw[Teal,dashed,line width=1.8pt] (2.1,1.7) rectangle (11.8,10.7);
        \foreach \x/\y in {
          .7/8.3,1.4/3.7,2.1/12.5,2.8/6.1,3.5/1.0,4.2/9.8,4.9/4.4,
          5.6/13.1,6.3/7.2,7.0/2.4,7.7/11.4,8.4/5.2,9.1/13.7,9.8/8.8,
          10.5/3.1,11.2/10.2,11.9/6.5,12.6/1.8,13.3/12.0,1.0/10.9,
          3.1/8.0,5.0/1.9,6.8/12.6,8.8/9.9,10.8/5.8,12.9/7.7}
          \fill[NavyTwo] (\x,\y) circle (0.13);
        \draw[Amber,line width=3pt] (3.1,3.0) -- (10.9,10.8);
        \foreach \t in {0,...,12}{
          \pgfmathsetmacro{\u}{3.1+0.65*\t}
          \pgfmathsetmacro{\v}{3.0+0.65*\t}
          \fill[Amber] (\u,\v) circle (0.19);
          \draw[White,line width=.6pt] (\u,\v) circle (0.19);
        }
        \node[anchor=west,font=\fontsize{13}{15}\selectfont,text=Navy] at (0,14.6)
          {\(1+1\) cross-section schematic};
        \node[rotate=45,font=\bfseries\fontsize{12}{14}\selectfont,text=Amber] at (9.7,11.8)
          {planted chain};
      \end{tikzpicture}
      \end{center}
    \end{column}
    \begin{column}{0.49\linewidth}
      Start with a background \(\Pi_\rho\subset K\) satisfying F1--F2. Along a fixed timelike
      geodesic segment of length \(\tau_0\), add
      \tightdisplay{
        \Gamma_\rho,\qquad
        k_\rho=\Bigl\lceil A(m_d\rho)^{1/d}\tau_0\Bigr\rceil.
      }
      In every fixed dimension:
      \begin{itemize}
        \setlength\itemsep{2mm}
        \item \(\gamma_0\cap D\) is an interval;
        \item \(\operatorname{len}(\gamma_0\cap D)\le\tau(D)\);
        \item \(\operatorname{Vol}(D)\ge\kappa_d\tau(D)^d\).
      \end{itemize}
      The added mass fraction is only \(\Theta(\rho^{1/d-1})\), but it is perfectly aligned in the order.
    \end{column}
  \end{columns}
\end{posterbox}

\begin{scopebox}
  \textbf{Geometric realization used for the check:}
  \tightdisplay{
    M_\ell^d=(-\infty,0)_\eta\times\mathbb R^{d-1}_{\mathbf x},
    \qquad g=\frac{\ell^2}{\eta^2}(-d\eta^2+d\mathbf x^2)
  }
  on a precompact regional domain with a deep witness diamond.
\end{scopebox}

\begin{posterbox}{Dimension: theorem for every fixed \(d\ge2\)}
  In \(d\)-dimensional planar de Sitter, slice integration gives a fixed
  \(\kappa_{d,\beta}>0\) such that
  \tightdisplay{
    \operatorname{len}(\gamma_0\cap D)\le\tau(D)
    \le\kappa_{d,\beta}^{-1/d}\operatorname{Vol}(D)^{1/d}.
  }
  At Madsen's minimum scale, the normalized planted cost is
  \tightdisplay{
    O\!\left((\log\mathcal N_\rho)^{2/d-3/2}\right)\to0.
  }
  Thus the rates are \(-1/2\) in \(1+1\), \(-5/6\) in \(2+1\), and \(-1\) in
  \(3+1\). The geometric transfer and exact F3 normalization are proved, not conjectured.
\end{posterbox}

\end{column}

\hfill

% ------------------------------------------------------------------
% COLUMN 2
% ------------------------------------------------------------------
\begin{column}{0.345\textwidth}

\begin{posterbox}{Estimate 1: the planted chain is invisible to F2}
  Causal convexity makes \(\gamma_0\cap D\) an interval, and a uniform volume--trace comparison gives
  \tightdisplay{
    \#(\Gamma_\rho\cap D)
    \le A m_d^{1/d}\kappa_{d,\beta}^{-1/d}
       (\rho\operatorname{Vol}_g(D))^{1/d}+2.
  }
  With \(B_d:=A m_d^{1/d}\kappa_{d,\beta}^{-1/d}\), the additional F2 cost is
  \tightdisplay{
    \Xi_{A,d}(\rho)=
    B_d(\kappa_{d,\beta}c_*^{-d})^{1/d-1/2}
       (\log\mathcal N_\rho)^{2/d-3/2}
    +O((\log\mathcal N_\rho)^{-3/2})
    \longrightarrow0.
  }
  Hence the final configuration still satisfies F2 with a fixed constant \(K_d\).
\end{posterbox}

\begin{posterbox}{Estimate 2: the same chain is visible to F3}
  For the planted endpoints \(x_\rho\prec y_\rho\),
  \tightdisplay{
    \frac{\ell_{P_\rho}(x_\rho,y_\rho)}{(m_d\rho)^{1/d}}-\tau_0
    \ge(A-1)\tau_0-o(1).
  }
  Madsen's F3 bound would allow only
  \tightdisplay{
    C_d\!\left(\varepsilon^2+
      \frac{\log^{3/2}\mathcal N_\rho}{\mathcal N_\rho^{1/(2d)}}\right)\tau_0.
  }
  Choose \(A>1+C_d\varepsilon^2\). Then
  \(\mu=A-1-C_d\varepsilon^2>0\), and for large \(\rho\) the full F3 inequality fails by at least
  \(\mu\tau_0/3\).
\end{posterbox}

\begin{posterbox}{Why both statements coexist}
  \begin{center}
  \begin{tikzpicture}[x=1cm,y=1cm]
    \draw[-{Latex[length=5mm]},line width=2pt,Navy] (0,0) -- (16.5,0);
    \node[below,font=\fontsize{13}{15}\selectfont,text=Muted] at (15.8,-0.15) {\(\rho\to\infty\)};
    \fill[NavyTwo] (1.2,0) circle (.18);
    \fill[Amber] (8.2,0) circle (.21);
    \fill[Teal] (14.7,0) circle (.21);
    \node[above=3mm,align=center,font=\fontsize{14}{17}\selectfont,text=Navy] at (2.1,0)
      {background\\\(\Theta(\rho)\)};
    \node[above=3mm,align=center,font=\bfseries\fontsize{14}{17}\selectfont,text=Amber] at (8.2,0)
      {planted chain\\\(\Theta(\rho^{1/d})\)};
    \node[above=3mm,align=center,font=\fontsize{14}{17}\selectfont,text=Teal] at (14.2,0)
      {mass fraction\\\(\Theta(\rho^{1/d-1})\to0\)};
    \draw[Teal,line width=1.7pt,{Latex[length=4mm]}-] (4.3,-1.35) -- (1.3,-1.35);
    \draw[Teal,line width=1.7pt,-{Latex[length=4mm]}] (4.3,-1.35) -- (7.3,-1.35);
    \node[below=1mm,align=center,font=\fontsize{13}{16}\selectfont,text=Teal] at (4.3,-1.35)
      {normalized F2 cost\\\(O((\log\mathcal N_\rho)^{2/d-3/2})\to0\)};
    \draw[Red,line width=1.7pt,{Latex[length=4mm]}-] (11.3,-1.35) -- (8.4,-1.35);
    \draw[Red,line width=1.7pt,-{Latex[length=4mm]}] (11.3,-1.35) -- (14.2,-1.35);
    \node[below=1mm,align=center,font=\bfseries\fontsize{13}{16}\selectfont,text=Red] at (11.3,-1.35)
      {relative F3 bias\\\(\Theta(1)\)};
  \end{tikzpicture}
  \end{center}
  \vspace{2mm}
  The perturbation is asymptotically negligible for counts but macroscopically visible to chain height.
\end{posterbox}

\begin{scopebox}
  \centering
  \badge{NavyTwo}{O5: PROVED\_MODULO\_O3\_SCOPE}
  \hspace{6mm}
  \badge{Teal}{P5.2: PASS\_WITH\_SCOPE}
\end{scopebox}

\begin{posterbox}{What is already closed mathematically}
  \begin{itemize}
    \setlength\itemsep{2.5mm}
    \item \(d=2\): deterministic product-order construction and full geometric transfer;
    \item every fixed \(d\ge2\): de Sitter volume--trace lemma, fixed F2 constant and
          constant-margin F3 failure;
    \item exact \(V_3(\tau)\) and \(V_4(\tau)\) formulas checked symbolically;
    \item both standard curvature-norm readings pass;
    \item O3: explicit regional reading retained;
    \item O5: proved modulo that O3 scope.
  \end{itemize}
  No simulation, fitted threshold or asymptotic-rate-only claim is used.
\end{posterbox}

\end{column}

\hfill

% ------------------------------------------------------------------
% COLUMN 3
% ------------------------------------------------------------------
\begin{column}{0.319\textwidth}

\begin{posterbox}{What we ask you to adjudicate}
  \begin{enumerate}
    \setlength\itemsep{7mm}
    \item \key{Admissibility in the original causal-set sense.}
      Is an induced-order embedding with the stated number--volume control still “faithful” or
      “Planck-scale uniform” when \(\Theta(\rho^{1/d})\) deterministic points are aligned along one
      timelike geodesic?

    \item \key{Foundational prior art.}
      Do you know an earlier construction with the same configuration-level certificate:
      exact order + quantitative count uniformity + an explicit longest-chain/proper-time failure?
  \end{enumerate}

  \vspace{5mm}
  {\color{Muted}\textbf{Separate author query.} The exact regional scope intended by Madsen's
  Remark 5.4 is being put to Madsen, not to you.}
\end{posterbox}

\begin{posterbox}{Closest results we found}
  \renewcommand{\arraystretch}{1.24}
  \begin{tabular}{@{}>{\bfseries\color{Navy}}p{0.23\linewidth}p{0.71\linewidth}@{}}
    Müller & \textbf{Already has the dangerous conjunction:} his Theorem 2 mechanism leaves the
      \(d^-\) Hauptvermutung false even under “Planck-scale uniformness”, a count condition on causally convex
      subsets, while proper-time geometry can be macroscopically wrong. The exact separation is
      different: fixed \(K\), two geometries and nearby order laws; no fixed embedding with
      exact induced order and a chain-height/F3 certificate. His mechanism is a conformal
      deformation near a maximizer, not \(\Theta(\rho^{1/d})\) planted points.\\
    Henson & Changing a small number of causal relations can destroy embeddability;
      not the addition of aligned elements.\\
    Braun & “Order + number = geometry” from laws of every sample size, under \(d\ge3\);
      different quantifiers and object.\\
    Aghili et al. & Chain distributions for uniform sprinklings in several dimensions;
      no adversarial F2-without-F3 construction.\\
  \end{tabular}

  \vspace{5mm}
  \centering
  \badge{Amber}{MÜLLER: CLOSEST NEGATIVE PRECURSOR}
  \vspace{4mm}

  \badge{NavyTwo}{DOES NOT SUBSUME THIS CERTIFICATE}
\end{posterbox}

\begin{posterbox}{A reply can take one of three precise forms}
  \begin{tabular}{@{}>{\bfseries\color{Teal}}p{0.08\linewidth}p{0.86\linewidth}@{}}
    A. & Yes: under the stated reading, this is a genuine F1+F2-not-F3 counterexample in the
         original causal-set sense.\\[3mm]
    B. & No: faithful or Planck-scale-uniform embedding intends an exclusion; please identify it.\\[3mm]
    C. & Known: an earlier source already contains this same mechanism; please identify it.\\
  \end{tabular}
\end{posterbox}

\begin{scopebox}
  \warn{Two source-level qualifications remain explicit:}
  \vspace{3mm}
  \begin{center}
  {\ttfamily\bfseries\fontsize{16}{20}\selectfont
    CURVATURE\_NORM\_UNDEFINED / BOTH\_STANDARD\_READINGS\_PASS\\[2mm]
    REGIONAL\_SCOPE\_NOT\_LITERAL / EXPLICIT\_REGIONAL\_READING\_ADOPTED}
  \end{center}
\end{scopebox}

\begin{posterbox}[colback=Navy,colframe=Navy,coltitle=White]{Claim discipline}
  \color{White}
  No absolute priority claim is made. The mathematical claim is conditional on the two readings
  printed above. The requested expert judgment concerns interpretation and precedent, not an appeal
  to authority in place of proof.
\end{posterbox}

\end{column}
\end{columns}

\vfill

% Footer
\begin{tcolorbox}[
  enhanced,colback=NavyTwo,colframe=NavyTwo,boxrule=0pt,arc=0mm,
  left=10mm,right=10mm,top=5mm,bottom=5mm,width=\textwidth
]
{\fontsize{14.2}{18}\selectfont\color{White}
\textbf{Sources.}
N. Madsen, arXiv:2607.05840v1 (2026), Defs. 2.2, 2.6 and Remarks 4.7, 5.4;
O. Müller, arXiv:2503.01719v2 (2025), Thm. 2 and discussion on pp. 3--4;
J. Henson, gr-qc/0601121, §1.3;
M. Braun, \emph{Class. Quantum Grav.} 43 (2026), 045015;
G. W. Gibbons and S. N. Solodukhin, arXiv:0706.0603;
L. Bombelli, J. Lee, D. Meyer and R. Sorkin, \emph{Phys. Rev. Lett.} 59 (1987), 521--524.
\hfill
\textbf{Purpose:} one expert question · no priority claim.}
\end{tcolorbox}

\end{frame}
\end{document}
